\documentclass[12pt]{article}
\usepackage[margin=1in]{geometry}
\usepackage{hyperref}
\setlength{\parskip}{0.8\baselineskip}
\setlength{\parindent}{0pt}
\pagestyle{empty}

\begin{document}

\today

Dear Editor,

I am submitting our manuscript, ``Reliable Machine Learning Under Distribution Shift: Uncertainty-Aware Diabetes Risk Prediction from National Health Survey Data,'' for consideration as a research article in PLOS Digital Health.

This study examines whether the calibrated probabilities and conformal prediction guarantees used to make machine-learning risk models trustworthy actually hold up under real, non-synthetic distribution shift, using diabetes-risk prediction on the CDC's Behavioral Risk Factor Surveillance System as the setting. We test four independent shift dimensions, temporal (6-year), demographic subgroup, geographic (US Census region), and model family (deep ensemble vs.\ gradient-boosted trees), and find that shift magnitude, not merely its presence, determines whether these safeguards hold. Every comparison in the paper carries a 95\% bootstrap confidence interval.

We believe this work fits PLOS Digital Health's scope in machine learning and clinical decision support, particularly its emphasis on rigorous validation using real-world data. The manuscript uses only publicly available, de-identified survey data and involves no human-subjects research requiring IRB approval.

As the corresponding author is based in Ghana, a Research4Life Group A country, and this research received no external funding, we request the Research4Life publication fee waiver.

We confirm this manuscript is not under consideration elsewhere.

Sincerely,

Edmund Eric Gah

\end{document}
